\chapter{Results}

Testing cite \cite{AI-BOOK} and \space\coolcite{AI-BOOK}

% TODO: Input appropriate strings and integers in the table below
\begin{table}[H]
\centering
    \begin{tabular}{|l|r|r|r|r|r|r|}
    \hline
    \rowcolor[HTML]{EFEFEF} 
    \textbf{Instance name} &
      \multicolumn{1}{c|}{\cellcolor[HTML]{EFEFEF}\textbf{n}} &
      \multicolumn{1}{c|}{\cellcolor[HTML]{EFEFEF}\textbf{N}} &
      \multicolumn{1}{c|}{\cellcolor[HTML]{EFEFEF}\textbf{F}} &
      \multicolumn{1}{c|}{\cellcolor[HTML]{EFEFEF}\textbf{S}} &
      \multicolumn{1}{c|}{\cellcolor[HTML]{EFEFEF}\textbf{M}} &
      \multicolumn{1}{c|}{\cellcolor[HTML]{EFEFEF}\textbf{A}} \\ \hline
    \cellcolor[HTML]{EFEFEF}\textbf{string}     & int & int & int & int & int & int \\ \hline
    \cellcolor[HTML]{EFEFEF}\textbf{string}     & int & int & int & int & int & int \\ \hline
    \cellcolor[HTML]{EFEFEF}\textbf{string}     & int & int & int & int & int & int \\ \hline
    \end{tabular}
\end{table}

The columns N, F, S, M, and A are for the problems None, Few, Some, Many, and Alternate, respectively.

The table entries either give the answer, or contain `?' for those cases where we were unable to find a solution within reasonable time. We define reasonable time as being ...

For those questions where there is a reason for our inability to find a good algorithm, because the problem is hard, we wrote `?!'.

% TODO: Generate a (tab-seperated) text file 'results.txt'.
For the complete table of all results, see the tab-separated text file {\tt results.txt}.